\capitulo{7}{Conclusiones y trabajo futuro}

\section{Conclusiones}

En este Trabajo de Final de Grado se ha desarrollado un sistema para almacenar y analizar qué hace un usuario durante la creación 3D en Blender. Para ello se creó un controlador que registra la información necesaria para poder hacerlo sin perturbar al usuario mientras usa el programa.

Se constató durante el desarrollo del proyecto que Blender puede acceder a mucha información sobre las diferentes secuencias en él mismo. Por tanto, el sistema puede grabar todo lo referente a los objetos, los cambios de la geometría, los cambios en el modelado, y otros elementos significativos que ocurren mientras el usuario trabaja en él.

Uno de los aspectos más importantes de nuestro trabajo fue encontrar una manera de registrar todo esto de forma continua y estructurada. El controlador almacena los datos para poder analizarlos más fácilmente después. Esto permite examinar la evolución de la escena y cómo interactúa el usuario con las herramientas que Blender ofrece.

También se consiguió integrar el sistema directamente dentro de Blender mediante un \textit{addon}. Esto hace que el usuario pueda trabajar de forma normal mientras el sistema registra información en segundo plano. Gracias a esto, los datos recogidos se parecen más a situaciones reales de trabajo y no solo a pruebas preparadas en laboratorio.

Otro aspecto importante fue la detección de algunos eventos concretos durante el modelado. Por ejemplo, el sistema puede detectar cuándo se duplican objetos, cuándo se modifica geometría, cuándo se cambia de modo dentro de Blender o cuándo se realizan ciertos cambios sobre coordenadas UV. Toda esta información ayuda a entender mejor cómo trabaja el usuario durante el proceso de diseño.

Además, el sistema desarrollado permite guardar los datos dentro de los propios archivos \textit{.blend}. De esta forma, la información no se pierde cuando se cierra el programa y puede seguir utilizándose más adelante.

En general, este proyecto sirve como una base bastante útil para estudiar el proceso de modelado 3D y analizar cómo trabaja una persona dentro de Blender.

\section{Conclusiones técnicas}

A nivel técnico, el proyecto también permitió aprender muchas cosas sobre Blender y sobre el desarrollo de sistemas de captura de datos.

Se comprobó que la API de Blender permite acceder fácilmente a información interna del programa, como objetos, geometría o estados de la escena. Esto facilitó mucho el desarrollo del sistema.

También se utilizaron \textit{handlers} y temporizadores para detectar cambios mientras el usuario trabajaba en tiempo real. Gracias a esto, el sistema puede registrar eventos continuamente sin necesidad de que el usuario haga nada especial.

Otra parte importante fue el uso de comparaciones entre estados de la escena. Esto permitió evitar guardar información repetida y ayudó a reducir la cantidad de datos generados.

Durante el proyecto también se vio que guardar demasiada información puede afectar al rendimiento del programa. Por eso fue necesario buscar un equilibrio entre capturar suficientes datos y mantener Blender funcionando de forma fluida.

Además, algunos eventos complejos fueron más difíciles de detectar. Por ejemplo, ciertos cambios en geometría o en coordenadas UV necesitan técnicas más avanzadas porque no siempre basta con registrar comandos simples.

Todo esto demuestra que crear un sistema de captura de datos dentro de Blender es un trabajo bastante complejo y que requiere adaptarse al funcionamiento interno del programa.

\section{Limitaciones del sistema}

Aunque el sistema opera adecuadamente y cumple con los fines principales del proyecto, aún hay ciertas incompatibilidades.

El mayor problema surge porque el software ha sido diseñado exclusivamente para Blender; como tal, no funciona directamente en otros programas de modelado 3D.

También hay algunos eventos complejos que dependen mucho del estado interno del programa. Por eso, en ciertas situaciones pueden aparecer pequeñas imprecisiones en los datos recogidos.

Otra limitación es que en sesiones largas de trabajo se genera una gran cantidad de información. Esto hace necesario aplicar filtros y optimizaciones para evitar almacenar demasiados datos.

Además, el sistema de análisis todavía es bastante básico. Aunque ya permite estudiar información importante, todavía se podrían crear herramientas mucho más avanzadas para analizar mejor todos los datos recogidos.

\section{Líneas de trabajo futuras}

Este proyecto puede seguir creciendo en el futuro y ofrece muchas posibilidades de mejora.

Una de las mejores ideas sería mejorar el sistema de análisis usando técnicas avanzadas como minería de datos o inteligencia artificial. Esto le permitirá comprender mejor el comportamiento del usuario e identificar patrones en el modelo.

Además, se podría diseñar una nueva forma de presentar la información. Por ejemplo, mostrar gráficos, mapas o animaciones directamente en la visión 3D actual.

Se podría adaptar el sistema a otros programas de modelado 3D u otras aplicaciones de CAD. Así, podrían compararse distintas herramientas y diferentes procesos de diseño.

Así mismo, se debería seguir optimizando el rendimiento para minimizar la carga sobre Blender, sobre todo cuando se trabaja con escenas grandes o a lo largo de una sesión.

La inteligencia artificial puede integrarse directamente con la información recopilada. Esto podría hacer que los sistemas propongan acciones, predigan operaciones o ayuden a usuarios mientras modelan dentro de Blender.

Por último, también sería útil utilizar inteligencia artificial junto a los datos para construir sistemas capaces de ofrecer sugerencias. El objetivo es facilitar la enseñanza al analizar cómo aprenden los estudiantes a modelar en 3D y qué cambios observa durante el curso de la educación.

En resumen, este proyecto demuestra que es posible registrar y analizar automáticamente la interacción del usuario dentro de Blender. Además, abre muchas posibilidades para seguir investigando y desarrollar herramientas más inteligentes relacionadas con el modelado 3D.